\section{Ejecución de la búsqueda}
En esta sección se documenta la ejecución de la metodología de búsqueda definida. El proceso se desarrolló de forma estructurada, para cada motor o entorno de búsqueda se registraron las fuentes consultadas, los filtros aplicados y los resultados obtenidos.

\subsection{Aplicación de la cadena de búsqueda}

La ejecución y la revisión se distribuyeron entre distintos integrantes del equipo, manteniendo registros de las consultas y de los resultados. Las búsquedas se realizaron en diferentes motores y entornos, con el fin de obtener una base representativa del período 2022--2025.

En Google se identificaron soluciones que aplican IA generativa a tareas del ciclo de vida de requisitos. Se encontraron plataformas como StoriesOnBoard AI~\cite{storiesonboard}, Miro~\cite{miroai} y Beam AI~\cite{beamai}. También se identificaron herramientas integradas a entornos de trabajo ágil, como Copilot4DevOps~\cite{copilotfordevops1}, Modern Requirements4DevOps~\cite{modernreq1} y ClickUp~\cite{clickup}, orientadas a la generación y gestión de historias de usuario.

En GitHub la búsqueda se centró en repositorios activos de código abierto que emplean IA para generar o validar historias de usuario. Entre los repositorios identificados se encuentran melahn/user-story-generator, kapadias/story-craft~\cite{storycraft}, nilenso/storymachine~\cite{storymachine} y jonverrier/RequirementLinter~\cite{requirementlinter}.

En marketplaces de Atlassian y Azure se encontraron complementos basados en IA generativa para la gestión de proyectos, como Storygenie~\cite{storygenie}, AutoStory~\cite{autostory}, AI User Story y Test Case Generator for Jira. También se identificaron opciones que incorporan validaciones básicas en documentos de entrada o que agregan vinculación con fuentes dentro del ciclo de desarrollo.

\subsection{Uso de herramientas de inteligencia artificial para la búsqueda}

Se utilizó Perplexity~\cite{perplexity} en modo Investigation con un único prompt en inglés y con un rango temporal acotado (2022--2025). Se restringieron los resultados a herramientas reales y utilizables, tanto open source como comerciales, excluyendo ideas sin implementación o pruebas de concepto. Para cada resultado se registraron el nombre y el enlace oficial, el modo de uso (CLI, plugin o SaaS), la documentación disponible, la fecha o versión más reciente, la licencia y una breve descripción orientada a los aspectos de interés. Además, se limitó la salida a un conjunto breve, sin duplicados y con un máximo de diez resultados.

Con Gemini~\cite{gemini} se siguieron dos recorridos con la misma estructura, en tres mensajes consecutivos: listado inicial, refinamiento técnico y selección con pasos de verificación. La diferencia entre ambos recorridos fue el modo del primer mensaje (en el primer caso, Deep Research; en el segundo, modo normal).

En el primer recorrido, el inicio con Deep Research solicitó, en inglés, un listado de herramientas públicas entre 2022 y 2025 que cumplieran en conjunto los aspectos de interés definidos, evitando repeticiones. El resultado indicó que no se identificó una herramienta que integrara todos los aspectos de forma verificable en ese período. En iteraciones posteriores se mencionó a aqua ALM (AI Copilot) como candidata y se sugirió una verificación mínima en su versión SaaS, aunque quedó pendiente confirmarlo de manera concluyente.

En el segundo recorrido se repitió la misma secuencia, con el primer mensaje en modo normal. Se obtuvo un conjunto preliminar más amplio y luego se refinó la información con datos prácticos. El cierre incluyó una selección final con pasos para probar cada opción, incluyendo herramientas como GitHub Spec-Kit, RequirementLinter y AI Requirements Engineer Assistant, además de un prototipo como GeneUS, junto con indicaciones para su ejecución o prueba.

\subsubsection{Recuperación inicial de resultados}
Como resultado de la aplicación de la cadena de búsqueda y del uso de herramientas de inteligencia artificial, se obtuvo un conjunto consolidado de herramientas, plataformas y repositorios que cumplen parcial o totalmente con los aspectos de interés. La lista completa se incluye en el Capítulo \ref{cap:anexo}, Anexo 2, diferenciada por fuente de búsqueda y fecha de recuperación.

A este conjunto se le aplicará el procedimiento de inclusión y exclusión descrito en la metodología, con el fin de filtrar la base inicial y conformar el conjunto final de herramientas que se analizará en la siguiente etapa. Este procedimiento permitirá descartar resultados redundantes, no verificables o fuera del alcance definido, y asegurar que la muestra final sea coherente con los objetivos del estudio.

\section{Filtrado de resultados}

En esta sección se describe la obtención de la lista de las herramientas para el análisis. Se quiere filtrar todas las herramientas obtenidas con los criterios de inclusión/exclusión.  Se explica brevemente para cada herramienta si la misma fue filtrada o incluida en la lista.

\subsection{Primer filtrado}
% Las primeras 10 probadas por Romina
ScopeMaster \cite{scopemaster} menciona en su página que puede ser utilizado para el análisis de requisitos y diferentes funcionalidades de generación, como la de casos de uso. No menciona la generación de historias de usuario, ya que habla sobre pasos que vendrían después de tener ya las historias de usuario creadas, como creación de casos de prueba o diagramas.
Tiene una integración con Jira gratuita, llamada StoryAnalizer \cite{storyanalizer}
enfocada en analizar historias de usuario para la mejora, pero no clarifica la generación automática de ellas. Por lo que es filtrada.

POPal, AI Test Case \& User Story Generator for Jira \cite{popal},  en su página tiene un video explicativo donde muestra que se puede generar historias de usuario a partir de épicas, por lo que pasa el filtrado. Y tiene prueba gratuita por lo que puede ser probada en el análisis.

StoryCraft \cite{storycraft},  herramienta con repositorio disponible, siendo su propósito la generación de historias de usuario por lo que es incluida.

StoryTest (by ScopeMaster) \cite{storytest}, descartada dado que el sitio no permite ser probado a la fecha y que la herramienta está incluida en ScopeMaster.

Visual Narrator \cite{visualnarrator} descartada dado que la herramienta se utiliza para la visualización de historias de usuarios, no para la generación.

Featmap, \cite{featmap} repositorio archivado, no menciona generación de historias de usuario ni utilización de IA.

USeR (User Story eReviewer) \cite{user}, no tiene implementación disponible y similar a la anterior, se utiliza para validar la calidad de historias de usuario, no la generación.

Aqusa \cite{aqusa}, utilizada para la calidad de historias. Filtrada.

Reqsuite~\cite{reqsuite}, filtrada ya que la IA es utilizada para control de calidad y sugerencias, no para la generación de historias de usuario.

Por lo tanto, de las diez primeras herramientas, quedaron dos para probar directamente el producto: POPal y StoryCraft.
% Hasta aca lo probado por Romina

% Le tocaba a Valentin pero Probado por Romina
Modern Requirements4DevOps \cite{modernreq1} menciona en su página \cite{modernreq2} que la herramienta puede ser utilizada para generar historias de usuario. Copilot4DevOps \cite{copilotfordevops1} es un módulo incluido dentro de Modern Requirements4DevOps y permite la generación de historias de usuario según su blog \cite{copilotfordevops2}. Dado que Copilot4DevOps forma parte de Modern Requirements4DevOps, ambas se consideran como una única solución a los efectos de este análisis, aunque se registran por separado en las tablas de filtrado para documentar sus capacidades individuales.

Ambiguity Detection Using LLMs (ChaseMc-Clellan) \cite{ambiguitydect}, no presenta archivo de código para generación de historias de usuario. Similar RE-Ambiguity-Detection-Experiment (+OpenReq) \cite{reambiguity} descartado dado que es un experimento de detección de ambigüedad no generación de historias de usuario.

Reqnroll \cite{reqnroll} fue filtrada ya que no menciona la generación de historias de usuario, dado que consiste en una herramienta de pruebas.
% las primeras 5 de Valentin probado por Romina

Quality Modeller - Curiosity Software \cite{qualitymodeller} queda filtrada, ya que aunque la herramienta ofrece soporte para la automatización de pruebas y permite la importación de requisitos, no genera historias de usuario con criterios de aceptación en formato ágil.

aqua ALM (AI Copilot) \cite{aquaalm} no pasa el filtro, ya que si bien incorpora un asistente de inteligencia artificial orientado a acelerar tareas dentro del ciclo de vida de pruebas, no está diseñado para generar historias de usuario a partir de documentos de requisitos o especificaciones. Su enfoque principal se centra en la gestión del ciclo de vida de aplicaciones y la optimización de procesos de testing.

GitHub Spec-Kit \cite{spec-kit} no pasa el filtro, ya que su enfoque está centrado en trabajar con especificaciones formales y producir planes técnicos y listas de tareas, pero no genera historias de usuario ni criterios de aceptación a partir de documentos de requisitos.

AI Requirements Engineer Assistant \cite{ai-requirements-engineer-assistant} no pasa el filtro, ya que se orienta a analizar, validar y mejorar documentos de requisitos, detectando ambigüedades, extrayendo requisitos y respondiendo preguntas sobre el documento. La herramienta ofrece funcionalidades avanzadas de ingeniería de requisitos, pero no genera historias de usuario de forma automática, por lo que queda descartada.

GeneUS \cite{geneus}  queda filtrada, ya que es solo un artículo académico que describe una herramienta conceptual basada en GPT‑4 para generar historias de usuario y casos de prueba automáticamente, pero no hay evidencia pública de una herramienta funcional disponible.
% las ultimas 5 de Valentin probadas por Gonza


%%%%%%%%%%%%%%%%%%%%%%%%%%%%
AutoDev Editor \cite{autodeveditor} es un editor de contenido basado en IA distribuido como proyecto de código abierto. En su repositorio se lo describe como un editor sofisticado para la creación de distintos tipos de contenido, incluyendo explícitamente “blogs, articles, user stories, and more”, es decir, historias de usuario entre los formatos soportados. La herramienta integra capacidades de generación y edición asistida, utilizando modelos de lenguaje para producir y refinar texto dentro del propio editor. Dado que la generación de historias de usuario forma parte del caso de uso declarado, y que el proyecto se encuentra públicamente disponible y mantenido, AutoDev Editor pasa el Filtrado 1 y se incluye en el conjunto de herramientas a analizar.

Specifai \cite{specifai} se presenta como una plataforma impulsada por IA para gestionar el ciclo de vida de requisitos y documentación. Su repositorio oficial la describe como una solución que transforma la gestión de requisitos de proyecto, automatizando documentación, generación de tareas y otros artefactos mediante flujos guiados por IA. En el ecosistema asociado se detalla además un servidor MCP que permite interactuar con documentos de requisitos como business requirements, product requirements y user stories, lo que indica soporte explícito para historias de usuario como tipo de artefacto manejado por la herramienta. Por tratarse de una solución que toma requisitos como entrada y genera documentación estructurada, incluyendo user stories, Specifai cumple el criterio de generación explícita de historias y se incluye en el Filtrado 1.

StoryBot \cite{storybot} es una herramienta de línea de comandos (CLI) que integra modelos de lenguaje para generar historias de usuario. En su repositorio se la define de forma explícita como “AI Generated User Stories from your CLI”, es decir, un generador de historias de usuario desde la terminal. A partir de un prompt textual, la herramienta produce historias de usuario detalladas, pensadas para integrarse en el flujo de trabajo de desarrollo ágil. Dado que la funcionalidad principal documentada es precisamente la generación de historias de usuario, StoryBot se incluye en el conjunto de herramientas seleccionadas tras el Filtrado 1.

UST – User Story Tutor \cite{ust} es una herramienta propuesta en el ámbito académico cuyo propósito es ayudar a equipos ágiles a mejorar la calidad de historias de usuario ya escritas. La documentación del repositorio y el artículo asociado describen que UST recibe como entrada el texto de una historia de usuario en inglés y devuelve recomendaciones personalizadas de mejora, índices de legibilidad y una predicción de story points basada en modelos de machine learning. En ningún caso se plantea la generación de nuevas historias a partir de requisitos u otras fuentes; la herramienta actúa como tutor de calidad sobre historias existentes. Por tanto, UST no cumple el criterio de generación explícita de historias de usuario y se filtra en esta etapa.

UserStory QE Agent \cite{userstoryqe} es un complemento para Jira Cloud desarrollado por UST, orientado a la ingeniería de calidad de historias de usuario. Según su página en Atlassian Marketplace, la aplicación “drive quality and efficiency in Jira by detecting defects early in user stories”, extrayendo automáticamente la información de la historia desde Jira (descripciones, adjuntos, épicas) y creando campos personalizados con los resultados de la validación y el puntaje de calidad. Su foco está en detectar defectos y mejorar historias de usuario ya cargadas en Jira, no en generarlas desde cero a partir de documentos o especificaciones. En consecuencia, UserStory QE Agent se clasifica como herramienta filtrada en el Filtrado 1.

QAssert TestCase and UserStory Generator \cite{qassert} es una app para Jira que se define explícitamente como una herramienta “AI-powered” para la generación automática de casos de prueba y creación de historias de usuario. En su ficha de Atlassian Marketplace se indica que permite “generate AI-powered user stories and test cases inside Jira”, mostrando ejemplos donde se generan historias de usuario con estructura de rol, funcionalidad, beneficio y criterios de aceptación, y luego se derivan casos de prueba vinculados a partir de esos criterios. Dado que la generación de historias de usuario es una funcionalidad central claramente descrita, QAssert pasa el Filtrado 1 y se incluye en el conjunto de herramientas candidatas.

AI Agile Story/Bug Generation \cite{aiagilestory} es un complemento para Jira desarrollado por Aneto Software. En los listados públicos se la describe como una solución para “supercharge your Jira with AI-powered user stories and bug reports”, cuyo objetivo es asistir en la creación de historias de usuario y reportes de errores generados con IA, refinando los requisitos capturados en el sistema. Aunque la documentación disponible es más breve que en otras herramientas, el nombre del producto y su descripción enfatizan la generación de user stories como parte de su propuesta de valor. Por ello, AI Agile Story/Bug Generation se considera incluida tras el Filtrado 1. 

AI Requirements Copilot for Jira \cite{airequirementscopilot} es un asistente de requisitos para Jira que utiliza IA para analizar, interpretar y reescribir los requisitos de un proyecto. La documentación oficial indica que “revolucionizes this process by harnessing the power of AI to analyze, interpret, and rewrite your project requirements into the formats that best suit your team's methodologies, including Business Epics in SAFe, User Stories, and more”, es decir, permite transformar requisitos en diferentes formatos, entre ellos historias de usuario. Esta capacidad de reescritura automatizada desde requisitos hacia user stories cumple el criterio de generación de historias a partir de artefactos de entrada, por lo que la herramienta se incluye en el Filtrado 1.

Storygenie \cite{storygenie} es una aplicación basada en Atlassian Forge que apunta a “revolucionar” la planificación ágil. Su descripción en el marketplace señala que, impulsada por IA, “creates Epics and User Stories from your prompts, offering bulk or individual generation”, permitiendo generar épicas e historias de usuario a partir de descripciones proporcionadas por el usuario, ya sea de forma masiva o individual. Además, integra mecanismos para refinar los resultados y asociar las historias generadas con épicas en Jira. Dado que la creación de historias de usuario a partir de prompts es el objetivo principal documentado, Storygenie se incluye en el conjunto de herramientas seleccionadas tras el Filtrado 1.

AutoStory: AI-Powered Ticket Creation for Jira \cite{autostory} es un complemento para Jira cuyo propósito es automatizar la creación de tickets mediante ChatGPT. La página oficial describe que la app “streamline your Jira ticketing process with ChatGPT AI-generated user stories and bug tickets”, generando historias de usuario y tickets de bug completos a partir de un título y una breve descripción, incluyendo descripción detallada, criterios de aceptación y detalles de prueba. La generación automática de user stories se presenta como beneficio central del producto, por lo que AutoStory cumple claramente el criterio de generación explícita de historias de usuario y se incluye en el Filtrado 1.

QVscribe \cite{qvscribe} es una herramienta que analiza automáticamente la calidad de requisitos escritos en lenguaje natural. Puede ser útil en caso de querer dar un enfoque analítico a los documentos a procesar, pero queda descartada de este análisis por no cumplir con la condición del filtrado de generar historias de usuario.

StoriesOnBoard \cite{storiesonboard} es una herramienta que cumple con la generación de historias de usuario, por lo que pasa el filtro. Además, genera criterios de aceptación y facilita la integración de las historias en un “story map”, actuando como un gestor interno de historias dentro de la herramienta.

Beam AI \cite{beamai} pasa el filtro, ya que cumple con la funcionalidad de generar historias de usuario. La forma de generar estas historias es introduciendo “personas” y “objetivos específicos” para cada una, a partir de los cuales las historias son generadas, incluyendo sus respectivos criterios de aceptación. La herramienta elimina ambigüedades y asegura que las historias sean concisas, garantizando que no superen las 50 palabras.

ClickUp \cite{clickup} cumple con la función de generar historias de usuario. Permite adjuntar archivos y utilizar prompts para su creación, y también puede generar criterios de aceptación si se solicita. Al funcionar como un chat, es posible indicarle especificaciones particulares para las historias. Las historias generadas se representan como “tareas” dentro de la aplicación, lo que facilita su seguimiento mediante el gestor de proyectos incluido en la plataforma.

Quanter \cite{quanter} no cumple con la función de generar historias de usuario, por lo que no pasa el filtro. Esta herramienta permite analizar y mejorar historias existentes mediante IA, detectando ambigüedades, lenguaje técnico inadecuado y errores comunes, además de verificar la estructura de las historias. Permite también exportarlas a Jira o Azure DevOps.

Miro AI \cite{miroai} cumple con la función de generar historias de usuario utilizando inteligencia artificial. Utiliza el contenido de un tablero interno de la aplicación (notas, personas, ideas, feedback) como contexto para generar las historias de usuario, por lo que pasa el filtro.

RequirementLinter \cite{requirementlinter} no pasa el filtro, ya que no genera historias de usuario. Esta herramienta utiliza IA para revisar y mejorar historias de usuario o requisitos formales basándose en guías reconocidas de la industria, evalúa además la claridad, estructura y calidad de las mismas. Al centrarse únicamente en análisis y no en generación, queda descartada.

StoryMachine \cite{storymachine} es una herramienta de línea de comandos que genera historias de usuario a partir de documentos de requisitos de producto (PRD) y especificaciones técnicas. Utiliza inteligencia artificial para procesar estos archivos y producir las historias, incluyendo criterios de aceptación, por lo que pasa el filtro.

Agile-bot \cite{agilebot} pasa el filtro, ya que permite generar historias de usuario de manera automática a partir de documentos de requisitos, además de validarlas mediante modelos de IA para asegurar coherencia y alineación con los objetivos del proyecto que se esté desarrollando. Ofrece funcionalidades de planificación de sprint. Su enfoque combina procesamiento con LLMs, embeddings y agentes para optimizar un flujo de trabajo dentro de un proyecto.

User Story Generator \cite{userstorygenerator} no pasa el filtro ya que no utiliza IA para la generación de historias.

AI-Driven Agile Team \cite{aidrivenagileteam}, no pasa el filtro ya que consta de un ejemplo utilizando IA para generar historias para un caso en particular.

En la Tabla \ref{table:filtrado1} se listan todas las herramientas y su estado.

%Quality Modeller \
%%%%%%%%%%%%%%%%%%%%%%%%%%%%5

\begin{center}
\begin{longtable}{|p{6.5cm}|p{4cm}|p{1.8cm}|}
\caption{Filtrado a partir de documentación}
\label{table:filtrado1} \\
\hline
\textbf{Herramienta} & \textbf{Estado} & \textbf{Fecha de ejecución} \\
\hline
\endfirsthead

\caption{Filtrado a partir de documentación - Continuación} \\
\hline
\textbf{Herramienta} & \textbf{Estado} & \textbf{Fecha de ejecución} \\
\hline
\endhead

\hline
\endfoot

\hline
\endlastfoot

ScopeMaster / Story Analyser for Jira & Filtrada, no menciona generación de historias & 16/11/2025 \\
RequirementLinter & Filtrada, no genera historias de usuario & 16/11/2025 \\
POPal --- AI Test Case \& User Story Generator for Jira & \textbf{Incluida}, demo en página sobre generación & 16/11/2025 \\
StoryCraft & \textbf{Incluida}, menciona en su repositorio sobre generación & 16/11/2025 \\
StoryTest (by ScopeMaster) & Filtrada, prueba no disponible e incluida en otra herramienta & 16/11/2025 \\
Visual Narrator & Filtrada, herramienta para visualizar historias & 16/11/2025 \\
Featmap & Filtrada, repositorio archivado & 16/11/2025 \\
USeR (User Story eReviewer) & Filtrada, implementación no disponible y académico  & 16/11/2025 \\
reqSuite RM & Filtrada, no se enfoca en la generación de historias & 16/11/2025 \\
AQUSA (Automatic Quality User Story Artisan) & Filtrada, calidad de historias no generación & 16/11/2025 \\
Modern Requirements4DevOps & \textbf{Incluida}, contiene Copilot4DevOps & 18/11/2025 \\
Copilot4DevOps & \textbf{Incluida}, genera historias de usuario con IA & 18/11/2025 \\
Ambiguity Detection Using LLMs (ChaseMcClellan) & Filtrada, el archivo sobre generación de historias de usuario está vacío & 18/11/2025 \\
RE-Ambiguity-Detection-Experiment (+ OpenReq) & Filtrada, experimento sobre ambigüedad & 18/11/2025 \\
Reqnroll & Filtrada, se utiliza para pruebas & 18/11/2025  \\
Curiosity Software - Quality Modeller & Filtrada, ofrece automatización de pruebas y soporte de requisitos, pero no genera historias de usuario de manera automática & 18/11/2025 \\
aqua ALM (AI Copilot) & Filtrada, asistente de IA para acelerar tareas de testing, no genera historias de usuario & 18/11/2025 \\
GitHub Spec-Kit & Filtrada, trabaja con especificaciones y planes técnicos, no genera historias de usuario ni criterios de aceptación & 18/11/2025 \\
AI Requirements Engineer Assistant & Filtrada, analiza y mejora requisitos, pero no genera historias de usuario automáticamente & 18/11/2025 \\
GeneUS & Filtrada, es un artículo académico sobre generación automática de historias de usuario, sin herramienta funcional disponible & 18/11/2025 \\
QVscribe & Filtrada, no genera historias de usuario & 15/11/2025 \\
Automated User Story Generation with Test Case Specification (GPT-4.0 based) & Filtrada, paper sin herramienta funcional disponible & 15/11/2025 \\
StoriesOnBoard & \textbf{Incluida}, genera historias y criterios de aceptación & 15/11/2025 \\
Beam AI User Stories & \textbf{Incluida}, genera historias con criterios de aceptación & 15/11/2025 \\
ClickUp & \textbf{Incluida}, genera historias mediante prompts y adjuntos, además con criterios de aceptación si se especifica & 15/11/2025 \\
Quanter AI & Filtrada, analiza y mejora historias, pero no genera & 15/11/2025 \\
Miro & \textbf{Incluida}, genera historias desde contenido del tablero interno de la aplicación & 15/11/2025 \\
RequirementLinter Public & Filtrada, revisa y mejora pero no genera historias & 15/11/2025 \\
StoryMachine & \textbf{Incluida}, genera historias a partir de PRDs y especificaciones & 15/11/2025 \\
Agile-bot & \textbf{Incluida}, genera y valida historias automáticamente & 15/11/2025 \\
User Story Generator & Filtrada, no utiliza IA & 19/11/2025 \\
%Story-craft & - & - \\
AI-Driven Agile Team & Filtrada, es un ejemplo de prueba & 19/11/2025 \\
AutoDev Editor & \textbf{Incluida}, menciona generación de user stories en su documentación & 16/11/2025 \\
Specifai & \textbf{Incluida}, plataforma de requisitos con generación de user stories & 16/11/2025 \\
StoryBot & \textbf{Incluida}, CLI explícitamente para generar user stories & 16/11/2025 \\
UST - User Story Tutor & Filtrada, tutor de calidad sobre historias existentes, no genera nuevas historias & 16/11/2025 \\
UserStory QE - Agent & Filtrada, valida y optimiza historias existentes, no se enfoca en su generación & 16/11/2025 \\
QAssert TestCase and UserStory Generator & \textbf{Incluida}, app de Jira para generación de user stories y casos de prueba con IA & 16/11/2025 \\
AI Agile Story/Bug Generation & \textbf{Incluida}, complemento de Jira para user stories y bugs generados con IA & 16/11/2025 \\
AI Requirements Copilot for Jira & \textbf{Incluida}, convierte requisitos en user stories y otros formatos ágiles & 16/11/2025 \\
Storygenie & \textbf{Incluida}, genera epics y user stories a partir de prompts & 16/11/2025 \\
AutoStory: AI-Powered Ticket Creation for Jira & \textbf{Incluida}, genera user stories y bugs en Jira usando ChatGPT & 16/11/2025 \\
\hline
\end{longtable}
\end{center}


Por lo tanto se tienen dieciocho herramientas que mencionan en su documentación la generación de historias de usuario. Estas se listan en la Tabla \ref{table:filtrado2}.

\begin{center}
\begin{longtable}{|p{7cm}|p{4cm}|p{1.8cm}|}
\caption{Herramientas Incluidas luego del filtrado general}
\label{table:filtrado2} \\
\hline
\textbf{Herramienta} & \textbf{Estado} & \textbf{Fecha de ejecución} \\
\hline
\endfirsthead

\caption{Herramientas Incluidas luego del filtrado general - Continuación} \\
\hline
\textbf{Herramienta} & \textbf{Estado} & \textbf{Fecha de ejecución} \\
\hline
\endhead

\hline
\endfoot

\hline
\endlastfoot

POPal --- AI Test Case \& User Story Generator for Jira & \textbf{Incluida}, demo en página sobre generación & 16/11/2025 \\
StoryCraft & \textbf{Incluida}, menciona en su repositorio sobre generación & 16/11/2025 \\
Modern Requirements4DevOps & \textbf{Incluida}, contiene Copilot4DevOps & 18/11/2025 \\
Copilot4DevOps & \textbf{Incluida}, genera historias de usuario con IA & 18/11/2025 \\
StoriesOnBoard & \textbf{Incluida}, genera historias y criterios de aceptación & 15/11/2025 \\
Beam AI User Stories & \textbf{Incluida}, genera historias con criterios de aceptación & 15/11/2025 \\
ClickUp & \textbf{Incluida}, genera historias mediante prompts y adjuntos, además con criterios de aceptación si se especifica & 15/11/2025 \\
Miro & \textbf{Incluida}, genera historias desde contenido del tablero interno de la aplicación & 15/11/2025 \\
StoryMachine & \textbf{Incluida}, genera historias a partir de PRDs y especificaciones & 15/11/2025 \\
Agile-bot & \textbf{Incluida}, genera y valida historias automáticamente & 15/11/2025 \\
AutoDev Editor & \textbf{Incluida}, menciona generación de user stories en su documentación & 16/11/2025 \\
Specifai & \textbf{Incluida}, plataforma de requisitos con generación de user stories & 16/11/2025 \\
StoryBot & \textbf{Incluida}, CLI explícitamente para generar user stories & 16/11/2025 \\
QAssert TestCase and UserStory Generator & \textbf{Incluida}, app de Jira para generación de user stories y casos de prueba con IA & 16/11/2025 \\
AI Agile Story/Bug Generation & \textbf{Incluida}, complemento de Jira para user stories y bugs generados con IA & 16/11/2025 \\
AI Requirements Copilot for Jira & \textbf{Incluida}, convierte requisitos en user stories y otros formatos ágiles & 16/11/2025 \\
Storygenie & \textbf{Incluida}, genera epics y user stories a partir de prompts & 16/11/2025 \\
AutoStory: AI-Powered Ticket Creation for Jira & \textbf{Incluida}, genera user stories y bugs en Jira usando ChatGPT & 16/11/2025 \\
\hline
\end{longtable}
\end{center}

\subsection{Segundo filtro}
Dadas las herramientas que pasaron el primer filtro, el siguiente paso es probar la generación de historias de usuario directamente en la plataforma.

En primera instancia fue posible descartar las herramientas que no se pudieron probar: para Beam AI User Stories no fue posible contactarse para conseguir acceso a la plataforma, AI Agile Story/Bug Generation dio error al intentar instalarla y AutoStory: AI-Powered Ticket Creation for Jira está desactualizada y no es posible probarla.

Para las herramientas que sí se obtuvo acceso se utilizó un documento técnico que contiene requisitos técnicos que permiten generar historias de usuario e incluye una inconsistencia que la herramienta podría detectar. Junto al documento también se utilizó, para las herramientas que lo requieran, un prompt simple que instruye al asistente a generar las historias, incluir criterios de aceptación, referencias al documento y conflictos.

\begin{verbatim}
    Analizar el documento adjunto/texto y generar historias
    de usuario a partir de los requerimientos del documento.
    Incluir criterios de aceptación para cada historia.
    Detectar y reportar inconsistencias encontradas en los
    documentos/texto de entrada, detectar cuando un requerimiento
    entra en conflicto con otro. Referenciar o vincular cada
    historia generada con una sección del documento de entrada.
\end{verbatim}

Para las herramientas que requieren prompting explícito, se utilizó un prompt simple como línea base: una instrucción breve que solicita las capacidades mínimas del estudio (generación de historias, criterios de aceptación, detección de conflictos y trazabilidad). En una segunda etapa se empleó un prompt completo alineado con la metodología del proyecto. El diseño y racional de ambos niveles de prompting se describen en la Sección~\ref{sec:diseno-prompts}.


De esta manera se comprobarán los requisitos: Generación de Historias con IA, Es automática (no requiere prompt), Inclusión de criterios de aceptación, Detección de inconsistencias, Permite Validar/Refinar las Historias, Trazabilidad.

POPal --- AI Test Case \& User Story Generator for Jira permite generar las historias, de manera automática sin requerir prompt e incluye criterios de aceptación. Además de permitir validar los resultados. Por lo tanto es incluida.

StoryCraft fue filtrada ya que requiere procesar requerimiento por requerimiento, no a partir del documento completo. % Valentin esto tiene sentido?

Modern Requirements4DevOps, la cual incluye Copilot4DevOps, fue incluida ya que contiene un chat que procesa el documento junto al prompt para generar historias que incluyen todos los requerimientos definidos. Copilot4DevOps tiene un botón automático para la generación, no requiere del prompt, aunque esto causa que se pierda el requisito de trazabilidad, pero el mismo puede ser editado.

StoriesOnBoard fue incluida en el filtrado inicial por su capacidad de generar
y refinar historias de usuario. Sin embargo, no permite la generación de
historias a partir de documentos de requisitos, sino que requiere la definición
manual de tarjetas que representan actores y sus funciones dentro del sistema, por lo que fue excluida del siguiente análisis.

ClickUp fue incluida, la misma cumple con todos los requisitos.

Miro fue incluida, requiere de prompt y en la prueba realizada generó las historias, incluyendo todos los requerimientos a excepción de la validación/refinado de historias, para la cual no se encontró la forma de realizarla en la primera prueba de la herramienta.

StoryMachine, fue incluida, es automática y no requiere de prompt, generó las historias a partir del documento, incluyendo criterios de aceptación y validación de las mismas.

Agile-bot fue filtrada, ya que no está orientada a historias de usuario sino a tareas técnicas. No detecta inconsistencias. La trazabilidad es básica pero útil. No permite refinar ni mejorar historias. No tiene criterios de aceptación.

AutoDev Editor fue incluida ya que genera historias de usuario. Permite elegir varios modelos para la generación, pero solo funciona uno y no se mantiene consistente en la generación cuando se refinan las historias. Por lo que no se puede refinar historias generadas, tampoco detectó las inconsistencias ni contiene trazabilidad.

Specifai fue incluida ya que genera en base a un título y una mínima descripción opcional, luego permite agregar requisitos en formato texto. Se puede iterar sobre lo generado, según la vista previa de la aplicación contiene un chat, pero no fue posible encontrarlo y utilizarlo en la prueba.

StoryBot fue filtrada ya que no fue posible la generación de las historias correctamente. % Keiver podrias ayudarme con la explicacion? No me queda claro

QAssert TestCase and UserStory Generator, fue incluida, la aplicación crea historias de usuario a partir de los requisitos de forma automática pero la salida solamente deja elegir de a una historia. Las mismas incluyen criterios de aceptación.

AI Requirements Copilot for Jira, fue incluida pero la generación de historias es a nivel de requerimiento, si se intenta hacer a nivel del documento de requerimientos no separa las historias en subhistorias automáticamente. Esto causa que no se pueda validar correctamente las historias generadas. Incluye criterios de aceptación.

Storygenie fue incluida, genera los registros con las historias de usuario, solo con un título y una descripción, no incluye criterios de aceptación, trazabilidad y tampoco detecta inconsistencias.

En la Tabla \ref{table:filtrado3}, se muestra el estado de cada herramienta luego de este filtrado.

\begin{center}
\begin{longtable}{|p{7cm}|p{4cm}|p{1.8cm}|}
\caption{Segundo Filtrado}
\label{table:filtrado3} \\
\hline
\textbf{Herramienta} & \textbf{Estado} & \textbf{Fecha de ejecución} \\
\hline
\endfirsthead

\multicolumn{3}{c}%
{{\bfseries \tablename\ \thetable{} -- continuación}} \\
\hline
\textbf{Herramienta} & \textbf{Estado} & \textbf{Fecha de ejecución} \\
\hline
\endhead

\hline \multicolumn{3}{r}{{Continúa en la siguiente página}} \\
\endfoot

\hline
\endlastfoot

POPal --- AI Test Case \& User Story Generator for Jira & \textbf{Incluida} & 21/11/2025 \\
StoryCraft & \textbf{Filtrada} & 23/11/2025 \\
Modern Requirements4DevOps & \textbf{Incluida} & 2/12/2025 \\
Copilot4DevOps & \textbf{Incluida} & 2/12/2025 \\
Beam AI User Stories & \textbf{Filtrada}, no se pudo probar & 23/11/2025 \\
ClickUp & \textbf{Incluida} & 13/10/2025 \\
Miro & \textbf{Incluida} & 27/11/2025 \\
StoryMachine & \textbf{Incluida} & 27/11/2025 \\
Agile-bot & \textbf{Filtrada}, no orientada a HU & 27/11/2025 \\
AutoDev Editor & \textbf{Incluida} & 01/12/2025 \\
Specifai & \textbf{Incluida} & 23/11/2025 \\
StoryBot & \textbf{Filtrada} & 25/11/2025 \\
QAssert TestCase and UserStory Generator & \textbf{Incluida} & 23/11/2025 \\
AI Agile Story/Bug Generation & \textbf{Filtrada}, no se pudo probar & 27/11/2025 \\
AI Requirements Copilot for Jira & \textbf{Incluida} & 29/11/2025 \\
Storygenie & \textbf{Incluida} & 21/11/2025 \\
AutoStory: AI-Powered Ticket Creation for Jira & \textbf{Filtrada}, no se pudo probar & 4/12/2025 \\
\hline
\end{longtable}
\end{center}

En la Tabla \ref{tab:incluidasfiltrado2} se listan todas las herramientas incluidas junto a los requisitos que cumplen.

\begin{longtable}{|p{2.5cm}|p{1.9cm}|p{1.2cm}|p{1.9cm}|p{2cm}|p{1.2cm}|p{1cm}|}
\caption{Herramientas incluidas luego del segundo filtrado} 
\label{tab:incluidasfiltrado2} \\
\hline
\textbf{Herramienta} &
\textbf{Generación automática de HU} &
\textbf{No utiliza prompt} &
\textbf{Incluye criterios de aceptación} &
\textbf{Detección de Inconsistencias} &
\textbf{Validar HU} &
\textbf{Traza- bilidad} \\
\hline
\endfirsthead

\caption{Herramientas incluidas luego del segundo filtrado - Continuación} \\
\hline
\textbf{Herramienta} &
\textbf{Generación automática de HU} &
\textbf{No utiliza prompt} &
\textbf{Incluye criterios de aceptación} &
\textbf{Detección de Inconsistencias} &
\textbf{Validar HU} &
\textbf{Trazabilidad} \\
\hline
\endhead

\hline
\endfoot

\hline
\endlastfoot

POPal --- AI Test Case \& User Story Generator for Jira & Sí & Sí & Sí & No & Sí & No \\
Modern Requirements4DevOps & Sí & No & Sí & Sí & Sí & Sí \\
Copilot4DevOps & Sí & Sí & Sí & Sí & Sí & No \\
ClickUp & Sí & No & Sí & Sí & Sí & Sí \\
Miro & Sí & No & Sí & Sí & No & Sí \\
StoryMachine & Sí & Sí & Sí & No & Sí & No \\
AutoDev Editor & Sí & No & Sí & No & No & No \\
Specifai & Sí & Sí & No & No & Sí & No \\
QAssert TestCase and UserStory Generator & Sí & Sí & Sí & No & No & No \\
AI Requirements Copilot for Jira & Sí & Sí & Sí & No & No & No \\
Storygenie & Sí & No & No & No & No & No \\
\end{longtable}

Por lo tanto tenemos once herramientas para evaluar sus historias de usuario (diez soluciones distintas, dado que Copilot4DevOps es un módulo de Modern Requirements4DevOps). Se decidió hacer un ranking dados los requisitos, para seleccionar las herramientas más completas.

\subsubsection{Ranking de los resultados}

Se armó un ranking con todas las herramientas que fueron probadas y que
efectivamente generan historias de usuario.

De acuerdo a la importancia que le demos a cada funcionalidad, definimos un valor
entre 0 y 1 para cada una, de forma tal que la suma total sea igual a 1.  
A este valor lo definimos como $W_i$.

Luego, para cada herramienta, definimos un valor $X_i$ que indica si la herramienta
cumple con la funcionalidad $i$.  
El valor de $X_i$ puede ser 0 o 1, según corresponda (en base a la tabla de evaluación).

Una vez definidos estos valores, el puntaje de cada herramienta se calcula de la siguiente manera:

\[
\text{puntaje\_herramienta} = \sum_{i} (X_i \cdot W_i)
\]

Es decir, para cada herramienta se inicializa el puntaje en 0 y, para cada funcionalidad $i$,
se suma el producto entre el cumplimiento de la funcionalidad y su peso asignado.
En la Tabla \ref{table:ranking-puntos} se presentan los valores definidos para cada funcionalidad según la importancia otorgada.

\begin{longtable}{|p{3cm}|c|p{3cm}|c|}
\caption{Peso de cada funcionalidad según importancia}
\label{table:ranking-puntos} \\
\hline
\textbf{Funcionalidad} & \textbf{Importancia} & \textbf{Justificación} & \textbf{Valor} \\
\hline
\endfirsthead

\caption{Peso de cada funcionalidad según importancia - Continuación} \\
\hline
\textbf{Funcionalidad} & \textbf{Importancia} & \textbf{Justificación} & \textbf{Valor} \\
\hline
\endhead

Trazabilidad & Menor &
Mínimo puntaje ya que muchas herramientas de generación no lo presentan. &
0,15 \\
\hline

Validación de historias generadas & Mayor &
Máximo puntaje, ya que es importante para el manejo y la calidad de las historias de usuario. &
0,25 \\
\hline

Detección de inconsistencias & Media &
Mediano puntaje, ya que esta funcionalidad es importante en el marco del proyecto.
Las inconsistencias en los documentos de requisitos pueden generar retrabajo en etapas
posteriores. Se considera más importante que la trazabilidad. &
0,20 \\
\hline

Inclusión de criterios de aceptación & Mayor &
Máximo puntaje, ya que determina la completitud de una historia de usuario. &
0,25 \\
\hline

No requiere prompt & Menor &
Mínimo puntaje, ya que la mayoría de las herramientas requieren un prompt. &
0,15 \\
\hline
\end{longtable}
